\begin{exercise}{15}
Show that for every $n\ge0$, there is a pushout
\[
 \begin{tikzcd}[column sep=large,row sep=large]
  \displaystyle\coprod_{J_n}\partial\Delta^n \arrow[r] \arrow[d]
    & \mathrm{sk}_{n-1}(X) \arrow[d] \\
  \displaystyle\coprod_{J_n}\Delta^n \arrow[r]
    & \mathrm{sk}_n(X)
 \end{tikzcd}
\]
\begingroup\sloppy
where $J_n$ is the set of non-degenerate $n$-simplices. Furthermore,
show that $X\cong\operatorname*{colim}_n\mathrm{sk}_n(X)$.
\par\endgroup

\begin{proof}
Put $\mathrm{sk}_{-1}(X)=\varnothing$.
By Lemma~1.1.39(1) and Exercise~\ref{ex:5}, the simplices of
$\mathrm{sk}_n(X)$ are exactly those whose unique nondegenerate
term in the decomposition of Exercise~\ref{ex:5} has dimension at most $n$.
For each $x\in J_n$, use the map $\Delta^n\to X$ corresponding to
$x$. A simplex $\alpha\colon[k]\to[n]$ in its boundary is not
surjective; factoring through its image as in Exercise~\ref{ex:3}
shows that $\alpha^*x$ lies in $\mathrm{sk}_{n-1}(X)$.
These maps define the displayed commutative square.

Check the pushout in degree $k$, using Proposition~1.1.23.
Its left vertical map is the inclusion of pairs $(x,\alpha)$ with
$x\in J_n$ and $\alpha\colon[k]\to[n]$ nonsurjective into the set
of all such pairs. Its complement consists of the pairs with
$\alpha$ surjective. By the uniqueness in Exercise~\ref{ex:5},
evaluation $(x,\alpha)\mapsto\alpha^*x$ bijects this complement
with
\[
 \mathrm{sk}_n(X)_k\setminus\mathrm{sk}_{n-1}(X)_k.
\]
Thus the pushout of sets identifies precisely the boundary pairs
with their evaluations in $\mathrm{sk}_{n-1}(X)_k$ and adds each
remaining simplex once. It is $\mathrm{sk}_n(X)_k$, as required.

Finally, $\mathrm{sk}_n(X)_k=X_k$ for $n\ge k$ by
Lemma~1.1.39(1). The degreewise colimit is therefore $X_k$ in every
degree. Proposition~1.1.23 now gives
$\operatorname*{colim}_n\mathrm{sk}_n(X)\cong X$.
\end{proof}
\end{exercise}
